\chapter{Business case, market research and objectives}
\label{ch:intro}

\section{The business problem}

Gym members rarely cancel on impulse. They stop coming first, and cancel weeks or months later; that lag is the window in which a well-chosen contact can still keep them. The window is valuable. The Health \& Fitness Association puts the median member retention rate at 66.4\% in 2024, across 175 operators running more than 17,000 facilities~\cite{hfa2025benchmarking}: about one member in three leaves each year, taking dues and ancillary spend with them. New members are the most fragile. In a study of one fitness centre, 63\% stopped before their third month and fewer than 4\% stayed continuously active beyond twelve~\cite{sperandei2016adherence}, and the real members analysed in this report agree: the newest segment churns at 94.6\% (\cref{sec:e1}).

The proposal identified three ways in which current practice fails (\cref{tab:practice}). The common thread is that each produces either a list or a score, but not a decision: none says whom to contact this week, with which offer, within a fixed budget and finite staff hours, and why.

\begin{table}[tb]
  \caption{Current retention practice and its failure modes, as identified in the proposal.}
  \label{tab:practice}
  \begin{tabularx}{\linewidth}{@{}l L@{}}
    \toprule
    \hd{Practice} & \hd{Why it fails} \\
    \midrule
    Calendar campaigns & Contact everyone on a schedule, so budget goes to members who would have stayed anyway. \\
    Threshold rules    & Ignore trend: twelve visits a month falling to four is a stronger warning than a steady two, but a threshold treats them alike. \\
    Vendor risk scores & Give a number, not a decision; staff cannot interrogate the flag, so they ignore it. \\
    \bottomrule
  \end{tabularx}
\end{table}

\section{Market and positioning}

Memberships grew 6\% and facilities nearly 4\% in 2024~\cite{hfa2025global}, so there are ever more members to keep. Club software divides into systems of record (membership, booking, payments, access control; for example ABC Fitness, Mindbody, Xplor, Perfect Gym, Virtuagym and Daxko) and an intelligence layer that bolts onto them, of which Keepme, with its predictive member score, is the clearest example; Glofox offers a built-in at-risk report. The proposal's desk research, based on vendor material and buyer guides that have a commercial interest, found four gaps:

\begin{enumerate}
  \item \textbf{Scores, not allocation.} No vendor publicly exposes constrained allocation: which flagged members to contact, with which offer, under a budget and staff hours.
  \item \textbf{Risk, not response.} Vendors target the members most likely to leave, but on field data, targeting by the expected response to the intervention reduces churn more than targeting by risk \cite{ascarza2018futility}.
  \item \textbf{Opaque scores.} Models trained on very large member pools are powerful but unauditable, and mid-market operators cite staff distrust of black-box flags.
  \item \textbf{Priced for chains.} The intelligence layer is generally recommended from five locations upwards, which leaves the one-to-five-club operator underserved.
\end{enumerate}

AnchorAI is therefore positioned as an \emph{explainable, budget-aware retention planner for mid-market operators}. It does not try to out-predict incumbents with far larger training data; it differentiates on the decision and on its transparency. We have not tested any competitor product, so every comparison in this report is against baseline policies in our own experiments, not against a vendor.

\section{Pricing and business value}
\label{sec:pricing}

% Pricing and business value, shared by both reports; each report writes its own heading before \input.
AnchorAI is sold as an annual subscription that bolts onto the operator's existing membership system: US\$199 a month for the first club and US\$149 for each further club up to five, so a three-club operator, the size of our scenario, pays US\$497 a month. It runs on the operator's own machine, language model included: in E8 a local seven-billion-parameter model was as faithful to its evidence as the hosted one (95.1\% of judged sentences supported for each), so there is no fee per question. A hosted and maintained option would cost more, to cover hosting and support; we have not costed it.

The price sits where the gap in the market is (\cref{tab:pricing}). The dedicated retention layer is priced for chains: in 2020 Keepme's founder quoted US\$9,000 a year for a single site, with a lower price for each further site~\cite{fittech2020keepme}, about US\$750 a month, nearly four times AnchorAI's first club. Small operators already pay at AnchorAI's level for their membership system and its add-ons~\cite{pushpress-pricing,gymdesk-pricing}.

\begin{table}[!htb]
  \caption{AnchorAI's price against published prices of club software, in US dollars a month.}
  \label{tab:pricing}
  \begin{tabularx}{\linewidth}{@{}l L r@{}}
    \toprule
    \hd{Product} & \hd{What it is} & \hd{Price a month} \\
    \midrule
    AnchorAI  & retention planner: first club; each further club, up to five & 199; 149 \\
    Keepme    & retention intelligence layer, single site (2020 figure)~\cite{fittech2020keepme} & about 750 \\
    PushPress & membership system; marketing add-on~\cite{pushpress-pricing} & 159--229; 329 \\
    Gymdesk   & membership system by member count; each extra location~\cite{gymdesk-pricing} & 75--240; 50--100 \\
    \bottomrule
  \end{tabularx}
\end{table}

The value comes from members kept. The planner values a kept member at six months of fees, a deliberately cautious horizon; at the scenario's middle fee of US\$54, for a six-month contract, that is US\$324, an assumption of the scenario layer rather than a market figure. One club therefore breaks even when AnchorAI keeps 0.6 members a month who would otherwise have left, and three clubs need 1.8. In the scenario, at the default 36 staff hours and US\$1,500, a week's plan from the greedy planner on the case-based estimates keeps 21.0 members against 13.7 for contacting by risk, and spends US\$432 instead of US\$1,280 (E4): 7.3 more members from a single week, four times the monthly break-even for three clubs. The scenario's effects are synthetic, so this shows that the margin is wide, not how wide it would be in a real gym. The direction agrees with field evidence that targeting by expected response reduces churn more than targeting by risk~\cite{ascarza2018futility}, and the stake is large, since about a third of members leave each year~\cite{hfa2025benchmarking}.


\section{Target user and approach}

The target user is a gym business with one to five clubs, and the people who use AnchorAI are its club managers, front-desk staff and trainers. No real operator was involved: we had no operator's data, and the manager interviews used to shape the rule base were simulated and are treated as plausible expert opinion, not evidence (\cref{sec:acquisition}).\footnote{The proposal named the team's study sponsor; it is not a client or data provider of this project and has not reviewed the system.}

AnchorAI answers the weekly question with one pipeline (\cref{ch:architecture}). It learns from real members who is at risk and why (\cref{ch:member}); checks the gym's playbooks and guardrails with a rule engine (\cref{ch:rules}); estimates what each offer would achieve for each member and chooses the best plan within the week's limits (\cref{ch:optimisation}); adds signals from a knowledge graph of who trains with whom; and lets staff question every recommendation in plain language (\cref{ch:knowledge}). Managers approve the plan, and every change is audited (\cref{ch:implementation}).
